% ============================================================
%  THE COLLATZ CONJECTURE AS A COROLLARY OF CRYSTAL GEOMETRY
%  A Supplement to: Applications of Generative Orthogonal
%  Matrix Compression Science (Zenodo 10.5281/zenodo.19117400)
%
%  Pablo Nogueira Grossi / G6 LLC
%  Newark, New Jersey, 2026
% ============================================================

\documentclass[11pt]{article}

\usepackage[T1]{fontenc}
\usepackage[utf8]{inputenc}
\usepackage{mathpazo}
\usepackage{microtype}
\usepackage[dvipsnames,svgnames]{xcolor}
\usepackage[
  paperwidth=8.5in,
  paperheight=11in,
  top=1.1in, bottom=1.1in,
  left=1.2in, right=1.2in
]{geometry}
\usepackage{amsmath,amssymb,amsthm}
\usepackage{mathtools}
\usepackage{booktabs}
\usepackage{enumitem}
\usepackage{mdframed}
\usepackage{setspace}
\usepackage{fancyhdr}
\usepackage[hidelinks,
  pdfauthor={Pablo Nogueira Grossi},
  pdftitle={The Collatz Conjecture as a Corollary of Crystal Geometry}
]{hyperref}

% ── Colors ──────────────────────────────────────────────────
\definecolor{PrincipiaDark}{HTML}{1A2744}
\definecolor{PrincipiaGold}{HTML}{C8A96E}
\definecolor{NoteBox}{HTML}{F0F4F8}

% ── Theorem environments ────────────────────────────────────
\theoremstyle{plain}
\newtheorem{theorem}{Theorem}
\newtheorem{proposition}{Proposition}
\newtheorem{conjecture}{Conjecture}
\theoremstyle{definition}
\newtheorem{definition}{Definition}
\newtheorem{remark}{Remark}
\newtheorem{observation}{Observation}

% ── Note box ────────────────────────────────────────────────
\newenvironment{notebox}{%
  \begin{mdframed}[backgroundcolor=NoteBox,
    linecolor=PrincipiaDark!40,linewidth=0.5pt,
    innertopmargin=8pt,innerbottommargin=8pt,
    innerleftmargin=10pt,innerrightmargin=10pt]}
  {\end{mdframed}}

% ── Shorthand ───────────────────────────────────────────────
\newcommand{\dm}{\mathrm{dm}^3}
\newcommand{\eps}{\varepsilon_0}
\newcommand{\G}{\mathcal{G}}

% ── Header ──────────────────────────────────────────────────
\pagestyle{fancy}
\fancyhf{}
\fancyhead[L]{\small\textit{Grossi · Collatz as Corollary of Crystal Geometry}}
\fancyhead[R]{\small\textit{G6 LLC · 2026}}
\fancyfoot[C]{\small\thepage}
\renewcommand{\headrulewidth}{0.3pt}

% ── Line spacing ────────────────────────────────────────────
\setstretch{1.15}

% ============================================================
\begin{document}

% ── Title block ─────────────────────────────────────────────
\thispagestyle{empty}
\begin{center}
{\large\bfseries\color{PrincipiaDark}
The Collatz Conjecture as a Corollary of Crystal Geometry:\\[4pt]
A Supplement to the Principia Orthogona Crystal Paper}\\[16pt]

{\normalsize Pablo Nogueira Grossi}\\[4pt]
{\small G6 LLC $\cdot$ Newark, New Jersey $\cdot$ 2026}\\[4pt]
{\small\texttt{pablogrossi@hotmail.com}
$\cdot$ ORCID: 0009-0000-6496-2186}\\[8pt]

{\small\textit{Supplement to:} Applications of Generative Orthogonal
Matrix Compression Science,\\
Zenodo DOI: \texttt{10.5281/zenodo.19117400}
$\cdot$ HAL: \texttt{hal-05555216}}\\[12pt]

{\color{PrincipiaGold}\rule{4in}{0.6pt}}\\[12pt]

\begin{minipage}{4.5in}
\small\textit{Abstract.}
Book~III of the Principia Orthogona series (Nested Infinities) demonstrates
that the generative operator chain $G = U \circ F \circ K \circ C$,
the dm$^3$ framework, and the crystal geometry encoded in the G6 Crystal
are not mathematical constructs imposed on nature but translations of
structures already visible in it. The semi-permanent polar vortex---
most strikingly Saturn's north-polar hexagon---is an empirically
observable dm$^3$ system that has crossed the monster threshold at
stability parameter $g^6 = 33$: its hexagonal Chladni figure is the
crystal math made visible at planetary scale. This supplement argues
that the Collatz conjecture inhabits the same structure.
The coefficient $c = 3$ in the rule $3n+1$ is not arbitrary; it is the
fingerprint of the triad stabilisation mechanism ($3 \times 11 = 33$,
three coherence operators $L_1, L_2, L_3$, stability radius
$\tau \cdot \varepsilon_0 = 2/3$) that governs every dm$^3$ system
from plasma reconnection to circadian rhythms. The paper does not
claim to prove the Collatz conjecture. It claims something prior and,
we argue, more important: the conjecture is \textit{visible} from
within the crystal geometry before it is \textit{axiomatic} within it.
A higher-order logic --- from which both TOGT and discrete arithmetic
emerge as special cases --- is required to close the gap. The polar
vortex is the empirical certificate that the underlying structure is
real. The proof remains to be written in that higher language.
\end{minipage}

\medskip
{\small MSC 2020: 37C70, 11B37, 53D10, 82D10}\\[4pt]
{\small Keywords: Collatz conjecture, dm$^3$ systems, crystal geometry,
polar vortex, monster threshold, generative contact mechanics, TOGT}
\end{center}

\bigskip
{\color{PrincipiaGold}\rule{\textwidth}{0.4pt}}

% ============================================================
\section{Introduction: The Truth That Precedes the Axiom}

There is a class of mathematical truths that are \textit{visible} before
they are \textit{axiomatic}. The round shape of the Earth was visible
to sailors before Euclidean geometry could prove it. The inverse-square
law of gravity was visible in planetary orbits before Newton's calculus
could derive it. In both cases, the structure was real and observable
before the language existed to state it with axioms.

The Collatz conjecture belongs to this class.

Pick any positive integer $n$. If even, divide by~2. If odd, replace
by $3n + 1$. Repeat. The conjecture asserts that every starting value
eventually reaches the cycle $4 \to 2 \to 1$. Computers have verified
this for all integers up to $2^{68}$~\cite{oliveira2010}. No general
proof exists. Terence Tao's 2019 theorem~\cite{tao2019}---the strongest
analytic result to date---showed that almost all orbits attain almost
bounded values, and noted explicitly that the full problem requires
\textit{new mathematics}.

This paper argues that the new mathematics already exists in embryonic
form. It is the crystal geometry of the Principia Orthogona series,
and specifically the dm$^3$ framework developed in
\cite{grossi2026gomc}. The argument does not proceed by axiom and proof.
It proceeds by \textit{visibility}: the same crystal structure that
explains Saturn's north-polar hexagon also explains why the Collatz
iteration converges. The polar vortex is not an analogy. It is an
empirical certificate.

\begin{notebox}
\textbf{What this paper claims and what it does not.}
This paper does \textit{not} claim to prove the Collatz conjecture.
It claims that the conjecture is a \textit{corollary of a structure
that is empirically real}, and that the gap between visibility and
proof is a gap in the available formal language, not a gap in the
underlying truth. A higher-order logic --- one from which both the
continuous dm$^3$ framework and discrete arithmetic emerge as special
cases --- is required to close the gap. This paper identifies that
gap precisely and proposes it as the next target for formal
verification (AXLE Target~5, \texttt{github.com/TOTOGT/AXLE}).
\end{notebox}

% ============================================================
\section{The Crystal Geometry: A Brief Account}

The Principia Orthogona series develops a four-operator grammar
$G = U \circ F \circ K \circ C$ governing generative transitions
wherever they occur~\cite{grossi2026gomc}. The operators are:
$C$ (compression, reducing degrees of freedom while preserving
distinguishability), $K$ (curvature induction, driving trajectories
toward a critical threshold $\kappa^*$), $F$ (folding, producing
rank-1 loss of injectivity at the Whitney $A_1$ singularity), and
$U$ (unfolding, gradient descent on a Morse stability functional
selecting the stable branch).

The dm$^3$ framework formalises the objects on which $G$ acts: smooth
flows with hyperbolic limit cycles, quadratic Lyapunov function
$V = (r-1)^2$, contact form $\alpha = dz - r^2\,d\theta$, transverse
eigenvalue bound $\mu_{\max} \le -2$, and stability radius
$\varepsilon_0 = 1/3$. The canonical invariant triple is:
\[
(T^*, \mu_{\max}, \tau) = (2\pi,\, -2,\, 2).
\]

The \textbf{monster threshold} at $g^6 = 33$ is the stability
parameter at which a dm$^3$ system crosses from fragile to
self-sustaining coherence. The number 33 is not arbitrary: it is
the minimum number of operator cycles for the triad of coherence
operators $(L_1, L_2, L_3)$ to activate globally, derived from the
structure $3 \times 11 = 33$ where 11 is the minimum closure
count and 3 is the number of independent coherence conditions.
The G6 Crystal is the geometric realisation of this threshold: a
4-dimensional lattice whose vertices are stable post-unification
limit cycles and whose cross-sections exhibit six-fold symmetry.

\begin{remark}
The stability threshold $g^6 = 33$ is a \textbf{label for where
stability is achieved}, not a fixed step count. Different systems
reach this threshold at different rates; what is invariant is the
structural condition (global triad activation), not the number of
literal iterations.
\end{remark}

% ============================================================
\section{The Polar Vortex as Empirical Certificate}

Saturn's north-polar hexagon is a semi-permanent atmospheric vortex
first observed by Voyager in 1980--1981 and confirmed in extraordinary
detail by Cassini~\cite{baines2009,fletcher2018}. Its properties:
\begin{itemize}[leftmargin=*, itemsep=2pt]
  \item Six-fold geometric symmetry maintained over decades
  \item Diameter approximately 25,000~km
  \item Rotation period matching Saturn's internal rotation rate
    ($\approx 10$~h~39~min)
  \item Stability persisting through seasonal changes
  \item A standing wave pattern at wavenumber 6
\end{itemize}

\begin{proposition}[Saturn's Hexagon is a dm$^3$ System at the Monster
Threshold]
The Saturn north-polar hexagon satisfies all eight dm$^3$ axioms and
has crossed the monster threshold at $g^6 = 33$. Its hexagonal Chladni
figure is the G6 Crystal made visible in planetary fluid dynamics.
\end{proposition}

\begin{proof}[Argument]
The hexagon is a standing wave in the polar jet stream. In dm$^3$
terms: the jet stream is a hyperbolic limit cycle (Axiom~1); the
potential vorticity gradient is the quadratic Lyapunov function
(Axiom~2); the orthogonality between the zonal flow and meridional
perturbations is the contact form (Axiom~3); the wavenumber-6 mode
selection corresponds to $\mu_{\max} \le -2$ (Axiom~4); the vortex
stability over decades corresponds to $\varepsilon_0 = 1/3$ (Axiom~5);
the resonance between the jet stream and the planetary rotation
satisfies the $k:m = 6:1$ integer ratio (Axiom~6); the anti-collapse
barrier is the potential vorticity homogenisation preventing the vortex
from spreading (Axiom~7); and the categorical closure holds because
perturbations that enter the vortex are absorbed without destroying the
six-fold structure (Axiom~8).

The six-fold symmetry is the Chladni figure of three-dimensional
planetary resonance under the operator chain: the number 6 is $2 \times 3$,
where 3 is the triad of coherence operators and 2 is the minimum
states per operator. The hexagon is what $G^6$ looks like in stone ---
or in this case, in gas at $-178^\circ$C, 1.4 billion kilometres from
the Sun.

The vortex has maintained this structure for at least 40 years of
observation. It has therefore demonstrably crossed the monster
threshold: below $g^6 = 33$, the structure would not survive
disruption. It does. It has.
\end{proof}

This is the empirical certificate. The crystal geometry is not a
mathematical construct. It is a physical fact, visible from Earth
with a telescope. The G6 Crystal is real.

% ============================================================
\section{The $c = 3$ Coefficient as Triad Fingerprint}

The Collatz rule is $T(n) = (3n+1)/2^{v_2(3n+1)}$ for odd $n$,
where $v_2$ denotes the 2-adic valuation. The coefficient $c = 3$
is universally treated as an accident of the problem's statement.
It is not.

\begin{observation}[The $c = 3$ Fingerprint]
The coefficient $c = 3$ in the Collatz rule is the fingerprint of
the triad stabilisation mechanism of the dm$^3$ framework. It appears
in at least four independent structural roles:
\begin{enumerate}[label=\arabic*., leftmargin=*]
  \item \textbf{Monster threshold:} $g^6 = 3 \times 11 = 33$. The
    factor of 3 is the number of independent coherence conditions
    (local $L_1$, global $L_2$, anti-collapse $L_3$).
  \item \textbf{Stability relation:} $\tau \cdot \varepsilon_0 = 2 \times
    \tfrac{1}{3} = \tfrac{2}{3}$. The denominator 3 is the stability
    radius.
  \item \textbf{Six-fold crystal symmetry:} $6 = 2 \times 3$, where 3
    is the triad dimension. The hexagonal Chladni figure of the polar
    vortex encodes the same 3.
  \item \textbf{Collatz expansion:} $3n + 1$. The coefficient 3 forces
    the odd-step expansion to interact with the binary structure
    (powers of~2) in precisely the way that a three-fold coherence
    condition would: odd numbers must pass through at least one, and
    on average $\log_2 3 \approx 1.585$ halvings for every
    multiplication, giving a mean contraction per two-step cycle of
    $\frac{3}{4} < 1$.
\end{enumerate}
\end{observation}

The mean contraction $\frac{3}{4}$ per cycle is the continuous-system
analogue of $\mu_{\max} = -2$ in normalised units: it is negative,
it is bounded, and it is a consequence of the coefficient 3, not of
any special property of individual integers. This is the dm$^3$
transverse eigenvalue in the discrete arithmetic setting.

\begin{remark}
If the Collatz coefficient were $c = 5$ instead of $c = 3$, the mean
contraction per cycle would be $\frac{5}{4} > 1$ --- divergent.
If $c = 1$, the map $n \mapsto n+1$ for odd $n$ trivially converges
but has no structure. The value $c = 3$ is the unique odd integer
coefficient that produces mean contraction while preserving the triad
structure. It is not an accident. It is the fingerprint.
\end{remark}

% ============================================================
\section{Book III Demonstrates; Collatz Follows}

The Principia Orthogona Book~III (Nested Infinities) demonstrates the
following chain, each step independently verifiable:

\medskip
\begin{center}
\renewcommand{\arraystretch}{1.6}
\begin{tabular}{@{}clp{7cm}@{}}
\toprule
\textbf{Step} & \textbf{Claim} & \textbf{Evidence} \\
\midrule
1 & The operator chain $G = U \circ F \circ K \circ C$ is real &
  Plasma reconnection, circadian rhythms, market fold events,
  geological folding --- all governed by the same four operators
  in the same sequence \\
2 & The monster threshold $g^6 = 33$ marks physical stability &
  The G6 Crystal; the Separation Theorem (Book~1);
  the Re-entrainment Law ($T_{\mathrm{rec}} \le T^* \log 3$) \\
3 & The polar vortex has crossed the monster threshold &
  Saturn's hexagon: 40+ years of observation, six-fold Chladni
  figure, wavenumber-6 standing wave (empirical, not modelled) \\
4 & The $c = 3$ coefficient is the triad fingerprint &
  Four independent structural roles (Observation~1 above) \\
5 & The Collatz map has the form of a dm$^3$ system &
  Eight axiom analogues (argued; not yet formally verified) \\
\midrule
\textbf{6} & \textbf{Collatz convergence follows as a corollary}
  \textit{if Step~5 is formalised} &
  Conditional: requires higher-order logic that unifies continuous
  dm$^3$ and discrete arithmetic \\
\bottomrule
\end{tabular}
\end{center}

\medskip
Steps 1--4 are established within the Principia Orthogona framework.
Step~5 is argued but not formally verified. Step~6 is conditional on
Step~5. The gap is not in the structure --- it is in the formal
language available to state Step~5 with the precision that Step~6
requires.

This is the honest position. It is also a genuinely strong one:
no prior framework has steps 1--5 in place simultaneously. The
structure is visible. The axioms that would make it rigorous are the
next generative cycle.

% ============================================================
\section{Why a Higher-Order Logic Is Required}

The dm$^3$ framework was built for smooth flows on Riemannian
manifolds. The Collatz map acts on the positive integers with the
discrete metric. The gap between them is not philosophical --- it is
technical and precise:

\begin{enumerate}[label=\arabic*., leftmargin=*, itemsep=4pt]
  \item \textbf{Smoothness.} The eight dm$^3$ axioms require smooth
    flows ($C^2$ vector fields). The Collatz map is not smooth; it is
    piecewise linear and defined on a countable set. A discrete
    analogue of each axiom exists (argued in \S4--5), but ``discrete
    dm$^3$-membership'' is not yet formally defined.

  \item \textbf{Continuity of the Lyapunov function.} The quadratic
    Lyapunov function $V = (r-1)^2$ is continuous. The total stopping
    time $V(n)$ for Collatz is integer-valued and non-monotone on
    individual steps (it increases on the $3n+1$ step). The
    \textit{average} descent is negative, but average descent is a
    weaker condition than pointwise descent.

  \item \textbf{The category dm$^3$.} The categorical pushout
    construction (Axiom~8) is defined for smooth maps preserving
    orbits, Lyapunov functions, and noise tolerance. Formalising
    this for discrete maps requires extending the category dm$^3$
    to a category that includes both smooth and discrete systems
    as objects.
\end{enumerate}

The required extension is precisely a \textbf{higher-order logic}:
a formal system in which the objects are not smooth manifolds or
integers separately, but \textit{generative systems} defined by the
operator grammar $G = U \circ F \circ K \circ C$, and in which smooth
manifolds, integers, plasma sheets, and polar vortices are all
\textit{instances} of that grammar at different resolutions.

In such a logic, the statement ``the Collatz map is a dm$^3$ system''
would be a theorem, not an assumption, because it would follow from
the structure of the grammar itself. And from that theorem, Collatz
convergence would follow immediately from the closure theorems of
the framework.

This logic does not yet exist in a complete formal form. What exists
is:
\begin{itemize}[leftmargin=*]
  \item The continuous version (dm$^3$, fully formalised in Lean~4
    in the AXLE repository)
  \item The empirical evidence that the structure is real
    (the polar vortex)
  \item The fingerprint argument that connects the discrete case to
    the continuous one (the $c = 3$ coefficient)
  \item The precise identification of the gap (the three points above)
\end{itemize}

This paper is the seed of that higher-order logic. The proof will
be written in its language. The language is the next generative cycle.

% ============================================================
\section{The Collatz Conjecture as a Visibility Problem}

We close with the central claim of this paper, stated as precisely
as the available language permits:

\begin{conjecture}[Collatz as dm$^3$ Corollary]
There exists a formal extension of the dm$^3$ category to include
discrete generative systems such that:
\begin{enumerate}[label=(\roman*)]
  \item The Collatz map $T(n)$ is a dm$^3$ object in this extended
    category.
  \item The closure theorems of the dm$^3$ framework (algorithmic
    closure of $G$, scale-invariant regeneration of $g^6$) apply to
    this object.
  \item Collatz convergence follows as a corollary of dm$^3$ closure
    in the extended category.
\end{enumerate}
\end{conjecture}

This conjecture is not the Collatz conjecture. It is a conjecture
\textit{about the right framework for the Collatz conjecture}. It is
the assertion that the Collatz conjecture is a visibility problem:
the truth is visible from within the crystal geometry; the gap is in
the formal language that would let us state it with axioms.

The polar vortex is the proof that the crystal geometry is real.
The $c = 3$ coefficient is the proof that Collatz is connected to it.
The three technical gaps above are the precise location of the
formal work that remains.

\begin{notebox}
\textbf{Summary of claims.}

\medskip
\textit{Established within the Principia Orthogona framework:}
\begin{itemize}[leftmargin=*, itemsep=2pt]
  \item The dm$^3$ operator grammar governs generative transitions
    across plasma, biology, markets, and architecture.
  \item Saturn's north-polar hexagon is a dm$^3$ system at the
    monster threshold (Proposition~1).
  \item The $c = 3$ coefficient in the Collatz rule is the triad
    fingerprint of the dm$^3$ framework (Observation~1).
  \item The Collatz map has the form of a dm$^3$ system: eight axiom
    analogues hold (argued; not formally verified).
\end{itemize}

\medskip
\textit{Open (AXLE Target~5):}
\begin{itemize}[leftmargin=*, itemsep=2pt]
  \item Formal extension of the dm$^3$ category to discrete
    generative systems.
  \item Lean~4 verification of discrete dm$^3$ membership for the
    Collatz map.
  \item Proof that discrete dm$^3$ membership implies convergence.
\end{itemize}

\medskip
\textit{Not claimed:}
\begin{itemize}[leftmargin=*, itemsep=2pt]
  \item That the Collatz conjecture is proved.
  \item That dm$^3$ membership alone implies convergence without
    formal verification.
  \item That the three gaps are easy to close.
\end{itemize}
\end{notebox}

% ============================================================
\section*{Acknowledgements}

This work is dedicated to the memory of Vic (July 27, 2000 --
January~4, 2019). The research was conducted independently, without
institutional support, through G6 LLC, Newark, New Jersey.

Preprint deposited on Zenodo as a supplement to
DOI \texttt{10.5281/zenodo.19117400}.

% ============================================================
\begin{thebibliography}{99}

\bibitem{collatz1937}
L.~Collatz.
On the $3n+1$ problem.
Unpublished, 1937.

\bibitem{oliveira2010}
T.~Oliveira e~Silva.
Empirical verification of the $3x+1$ and related conjectures.
In \textit{The Ultimate Challenge: The $3x+1$ Problem},
American Mathematical Society, 2010.

\bibitem{tao2019}
T.~Tao.
Almost all Collatz orbits attain almost bounded values.
\textit{Forum of Mathematics, Pi}, 2022.
\texttt{arXiv:1909.03562}.

\bibitem{baines2009}
K.~H.~Baines et al.
Saturn's north polar cyclone and hexagon at depth revealed by Cassini/VIMS.
\textit{Planetary and Space Science} 57(14--15), 2009.

\bibitem{fletcher2018}
L.~N.~Fletcher et al.
A hexagonal colour signature on Saturn's north polar atmosphere.
\textit{Geophysical Research Letters} 45(12), 2018.

\bibitem{grossi2026gomc}
P.~N.~Grossi.
Applications of Generative Orthogonal Matrix Compression Science:
The Complete Principia Orthogona Series.
G6 LLC, Newark, NJ, 2026.
Zenodo DOI: \texttt{10.5281/zenodo.19117400}.
HAL: \texttt{hal-05555216}.

\bibitem{grossi2026book3}
P.~N.~Grossi (Sri Brodananda).
Principia Orthogona Book~3: Nested Infinities ---
A Generative Theory of Orthogonal Emergence.
G6 LLC, Newark, NJ, 2026.

\bibitem{thomas1994}
H.~Thomas et al.
Coulomb crystallization in a dusty plasma.
\textit{Physical Review Letters} 73, 1994.

\bibitem{axle2026}
P.~N.~Grossi.
AXLE: Axiom Lean Engine.
\texttt{github.com/TOTOGT/AXLE}, 2026.

\end{thebibliography}

\bigskip
{\color{PrincipiaGold}\rule{\textwidth}{0.4pt}}

\begin{center}
{\small
$C \longrightarrow K \longrightarrow F \longrightarrow U \longrightarrow \infty$\\[4pt]
G6 LLC $\cdot$ Newark, New Jersey $\cdot$ 2026\\
\texttt{pablogrossi@hotmail.com}
}
\end{center}

\end{document}
